\icode{E-MM-SI-L3-261011-004}
如图，带等量异种电荷\(q\)的甲、乙两滑块，置于绝缘的水平桌面上，静止时相距为\(d\)。现在加一水平向右的匀强电场\(E\)，甲、乙两滑块开始运动，并发生碰撞结合在一起。已知甲的质量大于乙的质量，两者与桌面之间的动摩擦	因数相等，两滑块之间的相互作用不能忽略。下列说法正确的是\paren[]

\begin{tikzpicture}[font=\scriptsize]
	\draw[<->] (0,.25)--node[fill=white]{\(d\)}(2,.25);
	\foreach \x in {-1,-.8,...,2.8}
		\draw[gray] (\x,-.2)--++(.2,.2);
	\draw (-1,0)--(3,0);
	\foreach \y in {-.5,.1,.7,1.3}
		\draw[-latex,blue] (-1,\y)--(3,\y);
	\node at(-.25,.5)[above]{甲};
	\node at(2.25,.5)[above]{乙};
	\node at(3,.7)[above]{\(E\)};
    \draw[fill=red!20] (0,0) rectangle(-.5,.5);
    \draw[fill=blue!20] (2,0) rectangle(2.5,.5);
    \node at(-.25,.25)[red]{\(+ q\)};
	\node at(2.25,.25)[blue]{\(- q\)};
\end{tikzpicture}
\begin{choices}
	\item 甲、乙两滑块接近过程中任一时刻甲的速度比乙大
	\item 甲、乙两滑块相互接近过程中动量守恒
	\item 甲、乙发生碰撞结合在一起向左运动
	\item 从开始运动到碰撞的过程中，匀强电场和两电荷组成的系统的电势能减小量等于\(qEd\)	
\end{choices}
\begin{answer}
    C\qquad 提示：将甲乙看成一个系统，合外力向左，所以碰撞结合在一起向左运动
\end{answer}